\begin{statementbox}

	\textbf{\textcolor{CStatement}{条件}}
	\begin{description}[style=nextline, leftmargin=2.8em, labelsep=0.5em, itemsep=0.35em]
		\item[\textbf{\textcolor{CStatement}{条件 1（标准方程与焦点）：}}] 抛物线标准方程 $y^2 = 2px$，其中 $p>0$，焦点 $F(\frac{p}{2}, 0)$。
		\item[\textbf{\textcolor{CStatement}{条件 2（弦过焦点）：}}] 弦 $AB$ 经过焦点 $F$，且 $A,B$ 为抛物线上两个不同的点。
	\end{description}

	\textbf{\textcolor{CStatement}{结论}}
	\begin{description}[style=nextline, leftmargin=2.8em, labelsep=0.5em, itemsep=0.45em]
		\item[\textbf{\textcolor{CStatement}{结论一（倒数定值）：}}]
			\textbf{\textcolor{CStatement}{直观描述：}}
			焦点分弦所得两段长度 $|AF|$ 与 $|BF|$ 的倒数之和恒等于 $\frac{2}{p}$，与弦的倾斜角无关。
			\[
				\frac{1}{|AF|} + \frac{1}{|BF|} = \frac{2}{p}.
			\]

		\item[\textbf{\textcolor{CStatement}{推广结论（横坐标积为定值）：}}]
			\textbf{\textcolor{CStatement}{直观描述：}}
			设 $A(x_1,y_1)$，$B(x_2,y_2)$，则焦点弦两端点的横坐标之积为常数，与弦斜率无关。
			\[
				x_1 x_2 = \frac{p^2}{4}.
			\]
	\end{description}

\end{statementbox}
